%============================================================
%  ch03_algtensor.tex
%  Chapter 3: Algebraic Tensor Spaces
%
%  Tensor Formats in Banach Spaces
%  Antonio Falcó, CEU Universities
%
%  Estimated length: ~25 pp.
%============================================================

\chapter{Algebraic Tensor Spaces}
\label{chap:algtensor}

This chapter develops the purely algebraic foundations of tensor spaces.
We introduce the algebraic tensor product via its universal property,
establish the key notions of elementary tensor, matricisation, and multilinear
rank, and define the Tucker format and the set $\Tucker{\fr}{\VD}$ of tensors
of fixed Tucker rank. No topology is assumed in this chapter; that enters
in Chapter~\ref{chap:topTBT}.

Throughout, $D = \{1, \ldots, d\}$ with $d \geq 2$ is the \emph{dimension
index set}, and $V_\alpha$ ($\alpha \in D$) are vector spaces over $\KK$.
Concerning the foundational construction of tensor products we follow
Greub~\cite{Greub1981}.

%------------------------------------------------------------
\section{The algebraic tensor product}
\label{sec:algtensor:def}
%------------------------------------------------------------

\subsection*{Universal property and construction}

\begin{definition}
\label{def:algtensor}
The \emph{algebraic tensor product} of the spaces $V_\alpha$, $\alpha \in D$,
is a pair $(\VD, \bigotimes)$ consisting of a vector space $\VD$ and a
multilinear map
\begin{equation}
\label{eq:tensor_product_map}
    \bigotimes : \prod_{\alpha \in D} V_\alpha \longrightarrow \VD,
    \qquad
    (v_\alpha)_{\alpha \in D} \mapsto \bigotimes_{\alpha \in D} v_\alpha,
\end{equation}
satisfying the following \emph{universal property}: for every vector space
$Z$ and every multilinear map $M: \prod_{\alpha \in D} V_\alpha \to Z$,
there exists a unique linear map $\widehat{M}: \VD \to Z$ such that
$M = \widehat{M} \circ \bigotimes$.

We write $\VD = {}_a\bigotimes_{\alpha \in D} V_\alpha$ and refer to elements
of the form $\bigotimes_{\alpha \in D} v_\alpha$ (with $v_\alpha \in V_\alpha$)
as \emph{elementary tensors}. Every element of $\VD$ is a \emph{finite}
linear combination of elementary tensors.
\end{definition}

The algebraic tensor product exists and is unique up to isomorphism; see
\cite{Greub1981} for the standard construction. The prefix `$a$' in
${}_a\bigotimes$ emphasises the algebraic (non-topological) nature.
The underlying field is $\KK = \RR$ throughout, but all results hold
equally for $\KK = \CC$.

\begin{remark}
\label{rem:tensor_not_basis}
The set of elementary tensors is \emph{not} a basis of $\VD$; it
generates $\VD$ but with many linear dependencies. For example,
$\lambda v \otimes w = v \otimes \lambda w$ for all $\lambda \in \KK$.
If $\{e_i\}$ and $\{f_j\}$ are bases of $V_1$ and $V_2$ respectively,
then $\{e_i \otimes f_j\}$ is a basis of $V_1 \atens V_2$.
\end{remark}

\subsection*{Bilinear identification and the $\alpha$-unfolding}

For each $\alpha \in D$, the algebraic tensor space can be rewritten as
a two-factor product by grouping all indices except $\alpha$:
\begin{equation}
\label{eq:bilinear_split}
    \VD \cong V_\alpha \atens \Vbigalpha^{[\alpha]},
    \qquad
    \Vbigalpha^{[\alpha]} \coloneqq {}_a\bigotimes_{\beta \in D \setminus \{\alpha\}} V_\beta.
\end{equation}
Concretely, there is a canonical linear isomorphism
$\Phi_\alpha: \VD \xrightarrow{\;\sim\;} V_\alpha \atens \Vbigalpha^{[\alpha]}$
and throughout this monograph we identify $\bv \in \VD$ with
$\Phi_\alpha(\bv) \in V_\alpha \atens \Vbigalpha^{[\alpha]}$ without
further comment. An elementary tensor splits as
$\bigotimes_{\beta \in D} v_\beta \leftrightarrow v_\alpha \otimes
\bigotimes_{\beta \neq \alpha} v_\beta$.

\subsection*{Tensor products of linear maps}

Given linear maps $A_\alpha: V_\alpha \to W_\alpha$ for $\alpha \in D$,
their tensor product is the unique linear map
\begin{equation}
\label{eq:tensor_of_maps}
    \bigotimes_{\alpha \in D} A_\alpha :
    \VD \longrightarrow \mathbf{W}_D
    \coloneqq {}_a\bigotimes_{\alpha \in D} W_\alpha,
\end{equation}
determined on elementary tensors by
$\bigl(\bigotimes_{\alpha \in D} A_\alpha\bigr)
\bigl(\bigotimes_{\alpha \in D} v_\alpha\bigr)
= \bigotimes_{\alpha \in D} A_\alpha v_\alpha$.
This is the unique linear extension guaranteed by the universal property.
We denote the algebraic dual of $V$ by $\algdual{V} = L(V, \KK)$
and the space of all linear maps from $V$ to $W$ by $L(V, W)$.

\subsection*{Partitions and iterated tensor products}

\begin{definition}
\label{def:partition_tensor}
Let $\mathcal{P}$ be a partition of $D$ (i.e., a collection of
non-empty, pairwise disjoint subsets of $D$ whose union is $D$).
For each $\alpha \in \mathcal{P}$, set
$\Vbigalpha \coloneqq {}_a\bigotimes_{j \in \alpha} V_j$
(with the convention $\mathbf{V}_{\{j\}} = V_j$). Then the
\emph{iterated tensor product} associated to $\mathcal{P}$ is
\begin{equation}
\label{eq:iterated_tensor}
    {}_a\bigotimes_{\alpha \in \mathcal{P}} \Vbigalpha
    \cong \VD,
\end{equation}
where the isomorphism is canonical and consistent with the
identifications of elementary tensors.
\end{definition}

This associativity-of-tensor-products fact is fundamental: it allows us
to regroup the factors of $\VD$ according to any hierarchical structure,
which is precisely what the dimension partition trees of
Chapter~\ref{chap:trees} will exploit.

\begin{example}
\label{ex:iterated_tensor}
For $D = \{1,2,3,4\}$ and the partition
$\mathcal{P} = \{\{1,2\}, \{3,4\}\}$:
\begin{equation*}
    (V_1 \atens V_2) \atens (V_3 \atens V_4) \cong
    V_1 \atens V_2 \atens V_3 \atens V_4.
\end{equation*}
The partition $\mathcal{P}' = \{\{1\},\{2\},\{3\},\{4\}\}$ (trivial)
recovers the original four-fold product. The non-trivial partition
$\mathcal{P}$ is the structure underlying the Tucker format
(depth-2 tree), while more refined partitions correspond to hierarchical
formats; see Chapter~\ref{chap:trees}.
\end{example}

%------------------------------------------------------------
\section{Matricisation and multilinear rank}
\label{sec:algtensor:rank}
%------------------------------------------------------------

\subsection*{Matricisation in finite dimensions}

When the component spaces are finite-dimensional, the tensor $\bv \in \VD$
can be represented as a multi-dimensional array, and the identification
\eqref{eq:bilinear_split} induces a matrix.

\begin{definition}
\label{def:matricisation}
Let $D = \{1, \ldots, d\}$, $\dim V_\alpha = n_\alpha < \infty$, and
fix bases of each $V_\alpha$. Identify $\bv \in \VD$ with its
coordinate array $C \in \RR^{n_1 \times \cdots \times n_d}$.
For each $\alpha \in D$, the \emph{$\alpha$-matricisation} (or
\emph{$\alpha$-unfolding}) of $\bv$ is the matrix
\begin{equation*}
    \Mmat{\alpha}(\bv) \in \RR^{n_\alpha \times \prod_{\beta \neq \alpha} n_\beta}
\end{equation*}
obtained by reordering the coordinate array $C$ so that the $\alpha$-index
labels the rows and all other indices, in lexicographic order, label the
columns. The rank of this matrix,
$\rank \Mmat{\alpha}(\bv) \in \{0, 1, \ldots, n_\alpha\}$,
is the \emph{$\alpha$-rank} of $\bv$.
\end{definition}

The notion of $\alpha$-rank was introduced by Hitchcock~\cite{Hitchcock1927}
in 1927 under the name ``rank on the $j$th index''. In the
infinite-dimensional case, the correct replacement of
$\rank \Mmat{\alpha}(\bv)$ is the dimension of the minimal subspace
$\Umin{\alpha}{\bv}$, developed fully in Chapter~\ref{chap:minsubspaces}.

\subsection*{Explicit example: three-mode matricisation}

\begin{example}[Matricisation for $d=3$]
\label{ex:matricisation_d3}
Let $n_1=n_2=n_3=2$. Consider $\bv \in \RR^2 \otimes \RR^2 \otimes \RR^2$
with coordinate array
$C_{111}=C_{122}=C_{212}=C_{221}=1$, all other entries zero;
equivalently,
\begin{equation*}
    \bv = e_1{\otimes}e_1{\otimes}e_1
        + e_1{\otimes}e_2{\otimes}e_2
        + e_2{\otimes}e_1{\otimes}e_2
        + e_2{\otimes}e_2{\otimes}e_1.
\end{equation*}
The three matricisations, with columns ordered lexicographically by the
remaining indices, are:
\begin{align*}
    \Mmat{1}(\bv) &=
    \begin{pmatrix} 1 & 0 & 0 & 1 \\ 0 & 1 & 1 & 0 \end{pmatrix}
    \in \RR^{2 \times 4}, \\
    \Mmat{2}(\bv) &=
    \begin{pmatrix} 1 & 0 & 0 & 1 \\ 0 & 1 & 1 & 0 \end{pmatrix}
    \in \RR^{2 \times 4}, \\
    \Mmat{3}(\bv) &=
    \begin{pmatrix} 1 & 0 & 0 & 1 \\ 0 & 1 & 1 & 0 \end{pmatrix}
    \in \RR^{2 \times 4}.
\end{align*}
All three have rank~2, so the Tucker rank of $\bv$ is $(2,2,2)$ and
$\Umin{\alpha}{\bv} = \RR^2 = V_\alpha$ for each $\alpha$.
\end{example}

\subsection*{CP rank vs.\ Tucker rank}

\begin{example}[CP rank strictly exceeds Tucker rank]
\label{ex:cp_vs_tucker}
Consider the tensor $\bv$ of Example~\ref{ex:matricisation_d3}. Its
Tucker rank is $(2,2,2)$, yet its CP rank (the minimum number of
elementary summands) equals~4: the representation above has 4 terms,
and a 3-term representation can be ruled out by direct computation.

By contrast, the tensor
$\bw = e_1{\otimes}e_1{\otimes}e_1 + e_2{\otimes}e_2{\otimes}e_2
\in \RR^2 \otimes \RR^2 \otimes \RR^2$
has Tucker rank $(2,2,2)$ (each matricisation has rank~2) but CP rank~2.

This illustrates a fundamental distinction: the Tucker rank
$(r_1,\ldots,r_d)$ describes the subspace structure
($\Umin{\alpha}{\bv}$ has dimension $r_\alpha$), while the CP rank
counts minimal summands. One always has $r_\alpha \leq r_{\mathrm{CP}}(\bv)$,
but the inequality is strict in general for $d \geq 3$. Computing
Tucker rank is tractable (SVD of the matricisations); computing CP rank
is NP-hard~(H\aa stad, 1990). This computational advantage is a key
motivation for the Tucker-format framework of this monograph.
\end{example}

\begin{definition}
\label{def:fullrank_core}
For a rank tuple $\fr = (r_\alpha)_{\alpha \in D} \in \NN^d$, define
\begin{equation}
\label{eq:Rstar}
    \RR^{\times_{\alpha \in D} r_\alpha}_*
    \coloneqq
    \Bigl\{
        C \in \RR^{\times_{\alpha \in D} r_\alpha} :
        \rank \Mmat{\alpha}(C) = r_\alpha \;\text{ for all } \alpha \in D
    \Bigr\}.
\end{equation}
Elements of $\RR^{\times_{\alpha \in D} r_\alpha}_*$ are called
\emph{full-rank core tensors} of shape $\fr$.
\end{definition}

\begin{remark}
\label{rem:Rstar_open}
Since $\rank \Mmat{\alpha}(C) = r_\alpha$ if and only if
$\Mmat{\alpha}(C)\Mmat{\alpha}(C)^T \in \GL{\RR^{r_\alpha}}$,
and the determinant is a continuous function,
$\RR^{\times_{\alpha \in D} r_\alpha}_*$ is an open subset of
$\RR^{\times_{\alpha \in D} r_\alpha}$, hence a smooth finite-dimensional
manifold in its own right. In the scalar case $r_\alpha = 1$ for all
$\alpha$, we recover $\RR_* = \RR \setminus \{0\} = \GL{\RR}$.
\end{remark}

%------------------------------------------------------------
\section{The Tucker format}
\label{sec:algtensor:tucker}
%------------------------------------------------------------

\subsection*{Minimal subspaces: first appearance}

We introduce minimal subspaces here in their algebraic form; the full
theory (existence, uniqueness, lattice properties, best approximation)
is developed in Chapter~\ref{chap:minsubspaces}.

\begin{definition}
\label{def:minimal_subspace_algebraic}
Let $\bv \in \VD = {}_a\bigotimes_{\alpha \in D} V_\alpha$ and
$\alpha \in D$. Under the identification $\VD \cong V_\alpha \atens
\Vbigalpha^{[\alpha]}$, the \emph{minimal subspace} of $\bv$ in
direction $\alpha$ is the unique subspace $\Umin{\alpha}{\bv} \subset V_\alpha$
of smallest dimension such that
$\bv \in \Umin{\alpha}{\bv} \atens \Vbigalpha^{[\alpha]}$.
The \emph{$\alpha$-rank} of $\bv$ is $r_\alpha(\bv) \coloneqq \dim \Umin{\alpha}{\bv}$.
\end{definition}

The key properties — existence, uniqueness, and the symmetry
$\dim \Umin{\alpha}{\bv} = \dim \Umin{D \setminus \{\alpha\}}{\bv}$ —
are established in Chapter~\ref{chap:minsubspaces}. Here we proceed
taking these properties for granted.

\subsection*{The Tucker format and the set $\Tucker{\fr}{\VD}$}

\begin{definition}
\label{def:tucker_format}
Let $\fr = (r_\alpha)_{\alpha \in D} \in \NNz^d$ and let $U_\alpha \in
\Grass{r_\alpha}{V_\alpha}$ for each $\alpha \in D$. A tensor
$\bv \in \VD$ is said to be in \emph{Tucker format} with subspace
tuple $(U_\alpha)_{\alpha \in D}$ if
\begin{equation*}
    \bv \in {}_a\bigotimes_{\alpha \in D} U_\alpha.
\end{equation*}
The \emph{Tucker rank} of $\bv$ is the tuple
$\fr(\bv) = (r_\alpha(\bv))_{\alpha \in D}$,
and the set of tensors with fixed Tucker rank $\fr$ is
\begin{equation}
\label{eq:Tucker_set}
    \Tucker{\fr}{\VD} \coloneqq
    \bigl\{
        \bv \in \VD : r_\alpha(\bv) = r_\alpha \text{ for all } \alpha \in D
    \bigr\}.
\end{equation}
\end{definition}

The Tucker format thus represents $\bv$ as an element of a subspace
$U_1 \otimes \cdots \otimes U_d$ spanned by $r_\alpha$-dimensional
subspaces $U_\alpha \subset V_\alpha$. Unlike the subspaces, the
representation is not unique; the set $\Tucker{\fr}{\VD}$ collects all
tensors for which the \emph{minimal} such subspaces have dimension
exactly $r_\alpha$.

\subsection*{Admissible Tucker ranks}

\begin{definition}
\label{def:admissible_tucker}
A tuple $\fr = (r_\alpha)_{\alpha \in D} \in \NNz^d$ is an
\emph{admissible Tucker rank} for $\VD$ if there exists $\bv \in \VD$
with $r_\alpha(\bv) = r_\alpha$ for all $\alpha \in D$. The set of all
admissible Tucker ranks is
\begin{equation}
\label{eq:AD_tucker}
    \AD{\VD} \coloneqq
    \bigl\{
        (r_\alpha(\bv))_{\alpha \in D} \in \NNz^d : \bv \in \VD
    \bigr\}.
\end{equation}
\end{definition}

The zero tuple $\mathbf{0} = (0,\ldots,0)$ is always admissible
(corresponding to $\bv = 0$). The tuple $\mathbf{1} = (1,\ldots,1)$
is admissible if and only if $\dim V_\alpha \geq 1$ for all $\alpha$.
Since the $r_\alpha$ are bounded above by $\dim V_\alpha$ (in finite
dimensions), $\AD{\VD}$ is a finite set when all $V_\alpha$ are
finite-dimensional.

The following decomposition is one of the foundational structural facts
of the whole theory.

\begin{proposition}
\label{prop:tucker_decomposition}
The algebraic tensor space decomposes as a disjoint union:
\begin{equation}
\label{eq:tucker_disjoint}
    \VD = \bigsqcup_{\fr \in \AD{\VD}} \Tucker{\fr}{\VD}.
\end{equation}
\end{proposition}

\begin{proof}
Since every $\bv \in \VD$ has a well-defined Tucker rank tuple
(Chapter~\ref{chap:minsubspaces}), the sets $\Tucker{\fr}{\VD}$ cover
$\VD$; they are disjoint by definition of the rank.
\end{proof}

The decomposition \eqref{eq:tucker_disjoint} is the algebraic precursor
of the manifold decomposition established in
Chapter~\ref{chap:tucker_manifold}: each component $\Tucker{\fr}{\VD}$
will be shown to carry the structure of an infinite-dimensional
$\mathcal{C}^\infty$-Banach manifold.

\subsection*{Explicit representation of Tucker tensors}

The following lemma, which combines the definition of $\Tucker{\fr}{\VD}$
with the properties of minimal subspaces, gives three equivalent
characterisations of Tucker tensors. It is the algebraic heart of the
Tucker manifold construction.

\begin{lemma}
\label{lem:tucker_characterisation}
Let $\VD = {}_a\bigotimes_{\alpha \in D} V_\alpha$ and $\fr \in \AD{\VD}$
with $\fr \neq \mathbf{0}$. For $\bv \in \VD$, the following conditions
are equivalent.
\begin{enumerate}
\item[\textup{(a)}] $\bv \in \Tucker{\fr}{\VD}$.
\item[\textup{(b)}] \textup{(Tucker representation)} For each $\alpha \in D$ there exists a
    basis $\{u^{(\alpha)}_{i_\alpha} : 1 \leq i_\alpha \leq r_\alpha\}$
    of $\Umin{\alpha}{\bv}$ and a unique full-rank core tensor
    $\CD \in \RR^{\times_{\alpha \in D} r_\alpha}_*$,
    once the bases are fixed, such that
    \begin{equation}
    \label{eq:tucker_core}
        \bv = \sum_{\substack{1 \leq i_\alpha \leq r_\alpha \\ \alpha \in D}}
        \CD_{(i_\alpha)_{\alpha \in D}}
        \bigotimes_{\alpha \in D} u^{(\alpha)}_{i_\alpha}.
    \end{equation}
\item[\textup{(c)}] \textup{(Bilinear form)} For each $\alpha \in D$ there exist
    linearly independent vectors
    $\{u^{(\alpha)}_{i_\alpha}\}_{i_\alpha=1}^{r_\alpha}$ in $V_\alpha$
    and linearly independent vectors
    $\{\bU^{(\alpha)}_{i_\alpha}\}_{i_\alpha=1}^{r_\alpha}$
    in $\Vbigalpha^{[\alpha]}$ such that
    \begin{equation}
    \label{eq:tucker_bilinear}
        \bv = \sum_{i_\alpha=1}^{r_\alpha}
        u^{(\alpha)}_{i_\alpha} \otimes \bU^{(\alpha)}_{i_\alpha}.
    \end{equation}
    Moreover, when \eqref{eq:tucker_core} holds, the vectors
    $\bU^{(\alpha)}_{i_\alpha}$ are given by
    \begin{equation}
    \label{eq:tucker_Ualpha}
        \bU^{(\alpha)}_{i_\alpha} =
        \sum_{\substack{1 \leq i_\beta \leq r_\beta \\ \beta \in D \setminus \{\alpha\}}}
        \CD_{i_\alpha, (i_\beta)_{\beta \neq \alpha}}
        \bigotimes_{\beta \neq \alpha} u^{(\beta)}_{i_\beta}.
    \end{equation}
\end{enumerate}
\end{lemma}

\begin{proof}
\textbf{(a) $\Leftrightarrow$ (c).} If (a) holds, then by
Proposition~\ref{prop:min_d2} (Chapter~\ref{chap:minsubspaces}),
$\bv \in \Umin{\alpha}{\bv} \atens \Umin{D\setminus\{\alpha\}}{\bv}$
with $\dim \Umin{\alpha}{\bv} = \dim \Umin{D\setminus\{\alpha\}}{\bv} = r_\alpha$,
which gives~(c). Conversely, (c) implies $\dim \Umin{\alpha}{\bv} = r_\alpha$
for all $\alpha$, hence~(a).

\textbf{(b) $\Rightarrow$ (c).} Immediate.

\textbf{(c) $\Rightarrow$ (b).} From (c), one has
$\Umin{\alpha}{\bv} = \spann\{u^{(\alpha)}_{i_\alpha}\}$ for each $\alpha$,
and since $\bv \in {}_a\bigotimes_{\alpha \in D} \Umin{\alpha}{\bv}$,
there exists $\CD \in \RR^{\times_{\alpha \in D} r_\alpha}$ satisfying
\eqref{eq:tucker_core}. From \eqref{eq:tucker_Ualpha} and the linear
independence of the $\bU^{(\alpha)}_{i_\alpha}$, one deduces
$\rank \Mmat{\alpha}(\CD) = r_\alpha$ for all $\alpha$, so
$\CD \in \RR^{\times_{\alpha \in D} r_\alpha}_*$.
\end{proof}

\begin{remark}
\label{rem:tucker_identification}
From Lemma~\ref{lem:tucker_characterisation} and
Remark~\ref{rem:Rstar_open}, once bases of the minimal subspaces
$\Umin{\alpha}{\bv}$ are fixed ($\alpha \in D$), there is a natural
identification
\begin{equation}
\label{eq:tucker_iso}
    \Tucker{\fr}\Bigl({}_a\bigotimes_{\alpha \in D} \Umin{\alpha}{\bv}\Bigr)
    \xrightarrow{\;\sim\;}
    \RR^{\times_{\alpha \in D} r_\alpha}_*,
    \qquad
    \bu \mapsto E^{(D)},
\end{equation}
where $E^{(D)}$ is the coordinate array of $\bu$ in the chosen bases.
This identification is a diffeomorphism (both directions are polynomial),
and it is the local model for the Tucker manifold chart; see
Section~\ref{sec:tucker_manifold:charts}.
\end{remark}

\subsection*{The Tucker format as a flat case of tree-based formats}

The Tucker format corresponds to the \emph{shallowest} possible
hierarchical structure: the dimension partition tree $\TDtuck$ of
depth~1, whose root is $D$ and whose children are the singletons
$\{1\}, \ldots, \{d\}$:
\begin{equation*}
    \TDtuck = \bigl\{D, \{1\}, \{2\}, \ldots, \{d\}\bigr\},
    \qquad
    \Sons{D} = \bigl\{\{1\},\ldots,\{d\}\bigr\} = \Leaves[\TDtuck].
\end{equation*}
In Chapter~\ref{chap:algTBT} we will define tree-based formats for
arbitrary dimension partition trees $\TD$ and show that
$\Tucker{\fr}{\VD} = \FT{\fr}{\VD}{\TDtuck}$.

%------------------------------------------------------------
\section{The Hilbert--Schmidt norm as a tensor product norm}
\label{sec:algtensor:hs}
%------------------------------------------------------------

The Hilbert--Schmidt operators provide the most natural bridge between
the abstract tensor framework of this chapter and the concrete
infinite-dimensional setting of Part~\ref{part:quantum}. This section
shows that the Hilbert--Schmidt norm is a crossnorm on $H \otimes H^*$,
computes it explicitly for functions in $L^2$, and connects the
Hilbert--Schmidt rank with the Tucker rank.

\subsection*{The Hilbert tensor product}

Let $H$ be a separable complex Hilbert space with inner product
$\langle\cdot,\cdot\rangle$ and dual $H^*$ (Riesz isomorphism
$J: H \to H^*, J(v) = \langle v, \cdot\rangle$). The \emph{Hilbert
tensor norm} on $H \otimes_a H^*$ is defined for elementary tensors by
\begin{equation}
\label{eq:hilbert_tensor_norm}
    \|u \otimes \varphi\|_H \coloneqq \|u\|_H \|\varphi\|_{H^*},
\end{equation}
and extended to all finite sums by
\begin{equation}
\label{eq:hilbert_tensor_norm_sum}
    \Bigl\|\sum_k u_k \otimes \varphi_k\Bigr\|_H^2
    \coloneqq \sum_{j,k} \varphi_j(u_k)\,\overline{\varphi_k(u_j)}
    = \Tr\bigl(([\varphi_j(u_k)]_{jk})^* [\varphi_j(u_k)]_{jk}\bigr).
\end{equation}
The completion of $(H \otimes_a H^*, \|\cdot\|_H)$ is the space
$\cT{2}(H)$ of Hilbert--Schmidt operators, with
$\cT{2}(H) \cong H \otimes_H H^*$ isometrically
(Theorem~\ref{thm:app:hs_tensor} in Appendix~\ref{app:quantum}).

\begin{proposition}[Hilbert tensor norm is a crossnorm satisfying \condINJ]
\label{prop:hilbert_tensor_crossnorm}
\label{rem:grothendieck_hilbert}
The Hilbert tensor norm $\|\cdot\|_H$ on $H \otimes_a H^*$ satisfies:
\begin{enumerate}
\item[\textup{(i)}] \condCROSS: $\|u \otimes J(v)\|_H = \|u\|_H\|v\|_H$
    for all $u, v \in H$.
\item[\textup{(ii)}] \condINJ: $\|\cdot\|_\varepsilon \leq \|\cdot\|_H$.
\item[\textup{(iii)}] $\|\cdot\|_H \leq \|\cdot\|_\pi$ (dominated by
    the projective norm).
\item[\textup{(iv)}] The Grothendieck inequality gives
    $\|\cdot\|_\pi \leq K_G \|\cdot\|_\varepsilon$ on $H \otimes_a H^*$,
    so all three norms are equivalent with constant $K_G \approx 1.78$.
\end{enumerate}
\end{proposition}

\begin{proof}
Part~(i) is immediate from~\eqref{eq:hilbert_tensor_norm}.
Part~(ii): for any $\phi \in H^*$ with $\|\phi\|_{H^*} \leq 1$ and
$\psi \in (H^*)^*$ with $\|\psi\|_{(H^*)^*} \leq 1$, identifying
$\psi = J^{-1}(\psi) \in H$ (using $J$), one has
$|(\phi \otimes \psi)(\bv)| = |\phi((\text{Id} \otimes \psi)(\bv))|
\leq \|\bv\|_H$ by the Cauchy--Schwarz inequality in $\cT{2}(H)$,
giving $\|\bv\|_\varepsilon \leq \|\bv\|_H$.
Part~(iii) is standard: $\|\cdot\|_H \leq \|\cdot\|_\pi$ for any
reasonable crossnorm. Part~(iv) is Grothendieck's theorem
(Remark~\ref{rem:grothendieck_hilbert}).
\end{proof}

\subsection*{Hilbert--Schmidt rank equals Tucker rank}

\begin{proposition}[Hilbert--Schmidt rank and Tucker rank]
\label{prop:app:hs_rank}
Let $A \in \cT{2}(H)$ and identify $A$ with a tensor $\bA \in H \otimes_H H^*$.
Then $r_1(\bA) = r_2(\bA) = \mathrm{rank}(A)$, and the SVD
$A = \sum_{k=1}^r s_k v_k \otimes J(u_k)$ is the Tucker decomposition
with minimal subspaces $\Umin{1}{\bA} = \overline{\mathrm{Im}(A^*)}
\subset H$ and $\Umin{2}{\bA} = J(\overline{\mathrm{Im}(A)}) \subset H^*$.
\end{proposition}

\begin{proof}
Under the identification $\cT{2}(H) \cong H \otimes_H H^*$, the
contraction map $T_\varphi: H \otimes_H H^* \to H^*$,
$T_\varphi(\bA) = (\varphi \otimes \Id{H^*})(\bA) = \varphi(A(\cdot))$
has range $J(\mathrm{Im}(A))$ when $\varphi$ ranges over $H^*$. Hence
$\Umin{2}{\bA} = J(\overline{\mathrm{Im}(A)})$, and similarly
$\Umin{1}{\bA} = \overline{\mathrm{Im}(A^*)}$.
\end{proof}

\subsection*{Explicit computation in $L^2$}

\begin{example}[Hilbert--Schmidt norm of a separable kernel]
\label{ex:hs_l2}
Let $H = L^2([0,1])$ and $H^* \cong H$ via $J$. Consider the rank-$r$
kernel $K(x,y) = \sum_{k=1}^r \lambda_k \phi_k(x)\psi_k(y)$, viewed as
an element of $L^2([0,1]^2) \cong L^2([0,1]) \otimes_H L^2([0,1])$.
The associated Hilbert--Schmidt operator $T_K: L^2 \to L^2$,
$T_K f = \int_0^1 K(\cdot, y) f(y)\,dy$, has:
\begin{itemize}
\item \emph{Hilbert--Schmidt norm:}
    $\|T_K\|_{\cT{2}}^2 = \|K\|_{L^2([0,1]^2)}^2
    = \int_0^1 \int_0^1 |K(x,y)|^2\,dx\,dy$.
\item \emph{Tucker rank:} $r_1(K) = r_2(K) = r$
    (if $\{\phi_k\}$, $\{\psi_k\}$ are linearly independent).
\item \emph{Minimal subspaces:} $\Umin{1}{K} = \spann\{\phi_1,\ldots,\phi_r\}$
    and $\Umin{2}{K} = \spann\{\psi_1,\ldots,\psi_r\}$ in $L^2([0,1])$.
\end{itemize}
For the specific case $K(x,y) = xy$ (rank~1 kernel):
$\|K\|_{\cT{2}}^2 = \int_0^1 x^2\,dx \int_0^1 y^2\,dy = \tfrac{1}{9}$,
so $\|K\|_{\cT{2}} = \tfrac{1}{3}$. The Tucker rank is $(1,1)$ and
the Tucker decomposition is $K = \phi \otimes J(\psi)$ with
$\phi(x) = x$, $\psi(y) = y$.
\end{example}

\begin{remark}[The Hilbert tensor norm in the real case]
\label{rem:hs_real}
When $\KK = \RR$ and $H = \RR^n$, the Hilbert--Schmidt operators
on $\RR^n$ are exactly the $n \times n$ matrices with the Frobenius norm
$\|A\|_F = (\sum_{ij} A_{ij}^2)^{1/2}$. The identification
$\cT{2}(\RR^n) \cong \RR^n \otimes_H \RR^n$ (using $H^* \cong H$
since $J$ is the identity in the real self-dual case) gives
$\|A\|_F = \|A\|_H$, and the Tucker rank of $A$ viewed as a tensor in
$\RR^n \otimes \RR^n$ equals $\mathrm{rank}(A)$ as a matrix. This is
the finite-dimensional instance of Proposition~\ref{prop:app:hs_rank}.
\end{remark}

%------------------------------------------------------------
\section{The tensor product map and its differential}
\label{sec:algtensor:tensmap}
%------------------------------------------------------------

When the component spaces are normed, the tensor product map inherits
both topological and differential properties that underpin the manifold
constructions in Part~\ref{part:geometry}.

\begin{definition}
\label{def:tensor_product_continuous}
Let $(V_\alpha, \|\cdot\|_\alpha)$ be normed spaces and let $\|\cdot\|_D$
be a norm on $\VD$. The \emph{tensor product map}
\begin{equation}
\label{eq:tensor_product_normed}
    \bigotimes :
    \prod_{\alpha \in D}(V_\alpha, \|\cdot\|_\alpha) \longrightarrow
    (\VD, \|\cdot\|_D),
    \qquad
    (v_\alpha)_{\alpha \in D} \mapsto \bigotimes_{\alpha \in D} v_\alpha,
\end{equation}
(where the product is equipped with the $\max$ norm
$\|(v_\alpha)\|_\times = \max_\alpha \|v_\alpha\|_\alpha$) is said to be
\emph{continuous} if $\|\bigotimes_\alpha v_\alpha\|_D \leq C \prod_\alpha
\|v_\alpha\|_\alpha$ for some constant $C > 0$.
\end{definition}

Continuity of the tensor product map is equivalent to
$\|\bigotimes\| < \infty$; it holds whenever $\|\cdot\|_D$ is a
crossnorm (Definition~\ref{def:crossnorm}) on $\VD$.

\begin{proposition}
\label{prop:tensor_smooth}
Let $(V_\alpha, \|\cdot\|_\alpha)$ be normed spaces and let $\|\cdot\|_D$
be a norm on $\VD$ such that the tensor product map
\eqref{eq:tensor_product_normed} is continuous. Then:
\begin{enumerate}
\item[\textup{(i)}] The map \eqref{eq:tensor_product_normed} is
    $\mathcal{C}^\infty$-Fréchet differentiable, with differential
    \begin{equation}
    \label{eq:diff_tensor_product}
        D\bigotimes(v_1,\ldots,v_d)(w_1,\ldots,w_d)
        = \sum_{\alpha \in D}
          v_1 \otimes \cdots \otimes v_{\alpha-1} \otimes
          w_\alpha \otimes v_{\alpha+1} \otimes \cdots \otimes v_d.
    \end{equation}
\item[\textup{(ii)}] The tensor product of bounded linear maps,
    \begin{equation*}
        \bigotimes :
        \prod_{\alpha \in D} \cL(U_\alpha, V_\alpha)
        \longrightarrow
        \cL\Bigl({}_a\bigotimes_{\alpha \in D} U_\alpha,\, \VD\Bigr),
        \quad
        (A_\alpha) \mapsto \bigotimes_{\alpha \in D} A_\alpha,
    \end{equation*}
    is also continuous and $\mathcal{C}^\infty$-Fréchet differentiable
    for any finite-dimensional subspaces $U_\alpha \subset V_\alpha$.
\end{enumerate}
\end{proposition}

\begin{proof}
Part~(i) follows directly from Proposition~\ref{prop:multilinear_smooth},
since $\bigotimes$ is a continuous multilinear map. The formula for the
differential follows by the Leibniz rule for multilinear maps. For~(ii),
one checks boundedness of the map using the continuity of
\eqref{eq:tensor_product_normed}:
\begin{equation*}
    \|\bigotimes_\alpha A_\alpha\|_{\VD \leftarrow \bigotimes U_\alpha}
    \leq C \prod_\alpha \|A_\alpha\|_{V_\alpha \leftarrow U_\alpha},
\end{equation*}
where $C$ is the constant in Definition~\ref{def:tensor_product_continuous},
and the finite-dimensionality of $U_\alpha$ ensures
$L(U_\alpha, V_\alpha) = \cL(U_\alpha, V_\alpha)$.
Fréchet differentiability then follows again from
Proposition~\ref{prop:multilinear_smooth}.
\end{proof}

\begin{remark}
\label{rem:regularity_Tucker}
Proposition~\ref{prop:tensor_smooth} is the reason why the Tucker manifold
$\Tucker{\fr}{\VD}$ is $\mathcal{C}^\infty$ (not merely continuous): all
coordinate change maps between charts are built from compositions of
projections and the tensor product map, all of which are
$\mathcal{C}^\infty$-Fréchet differentiable. See
Section~\ref{sec:tucker_manifold:cinfty} for the precise argument.
\end{remark}

\subsection*{Banach tensor spaces}

When a norm is fixed on $\VD$, its completion is the natural ambient space.

\begin{definition}
\label{def:banach_tensor_space}
Let $(V_\alpha, \|\cdot\|_\alpha)$ be normed spaces and $\|\cdot\|_D$ a
norm on $\VD = {}_a\bigotimes_{\alpha \in D} V_\alpha$. The
\emph{Banach tensor space} associated to this data is the completion
\begin{equation*}
    \VDnorm
    \coloneqq
    {}_{\|\cdot\|_D}\bigotimes_{\alpha \in D} V_\alpha
    \coloneqq
    \overline{\VD}^{\|\cdot\|_D}.
\end{equation*}
If $\VDnorm$ is a Hilbert space, it is called a \emph{Hilbert tensor space}.
\end{definition}

\begin{example}
\label{ex:Lp_tensor}
For $I_\alpha \subset \RR$ ($\alpha \in D$) and $1 \leq p < \infty$,
the Lebesgue spaces $L^p(I_\alpha)$ satisfy
\begin{equation*}
    L^p\Bigl(\prod_{\alpha \in D} I_\alpha\Bigr)
    = {}_{\|\cdot\|_{0,p}}\bigotimes_{\alpha \in D} L^p(I_\alpha),
\end{equation*}
where $\|\cdot\|_{0,p}$ denotes the standard $L^p$ norm on the product
domain. For $p=2$ this is a Hilbert tensor space.
\end{example}

\begin{example}
\label{ex:sobolev_tensor}
For Sobolev spaces $H^{N,p}(I_\alpha)$ with $1 < p < \infty$ and
$\mathbf{I} = \prod_\alpha I_\alpha$:
\begin{equation*}
    H^{N,p}(\mathbf{I})
    = {}_{\|\cdot\|_{N,p}}\bigotimes_{\alpha \in D} H^{N,p}(I_\alpha).
\end{equation*}
This is reflexive and separable. For $p=2$ it is a Hilbert tensor space.
These spaces arise naturally in the numerical treatment of high-dimensional
PDEs; see Section~\ref{sec:min:bestapprox} and
Section~\ref{sec:topTBT:bestapprox}.
\end{example}

\begin{remark}
\label{rem:algebraic_dense}
The algebraic tensor space $\VD$ is dense in $\VDnorm$ by construction.
The Tucker manifold $\Tucker{\fr}{\VD}$ lives inside the algebraic space,
but its closure (as a subset of $\VDnorm$) and its topology as a manifold
are determined by the norm $\|\cdot\|_D$ and, in particular, by whether
$\|\cdot\|_D$ is a crossnorm. This interplay between algebraic and
topological structure is a recurring theme throughout the monograph.
\end{remark}

